% Faithful translation baseline; Chinese source PDF pp. 31--34.
\section{Schr\"odinger and Heisenberg Pictures}
\subsection*{1) Two kinds of unitary transformations}
The spatial-translation operator $\op U(\dd\vect x)$ and time-translation operator $\op U(\dd t)$ discussed earlier both act on the state:
\[\ket\alpha\longrightarrow\op U(\dd\vect x)\ket\alpha
=\left(\mathbf1-\frac{\ii}{\hbar}\vect{\op p}\cdot\dd\vect x\right)\ket\alpha,\]
\[\ket\alpha\longrightarrow\op U(\dd t)\ket\alpha
=\left(\mathbf1-\frac{\ii}{\hbar}\op H\dd t\right)\ket\alpha.\]
The physical quantities remain unchanged; observable differences arise from changes in the state. Under a unitary change of the state,
\[\mel\beta{\op A}\alpha\longrightarrow
\bra\beta\op U^{\dagger}\op A\op U\ket\alpha.\]
Equivalently, one may leave the state $\ket\alpha$ unchanged and transform the physical quantity:
\[\boxed{\op A\longrightarrow\op U^{\dagger}\op A\op U,\qquad\ket\alpha\ \text{unchanged}.}\]
The two methods are completely equivalent: the time dependence may be assigned either to the state or to the observable. For example, under a spatial displacement,
\[\op x\longrightarrow\op U^{\dagger}(\dd\vect x)\op x\op U(\dd\vect x)
=\left(\mathbf1+\frac{\ii}{\hbar}\vect{\op p}\cdot\dd\vect x\right)
\op x\left(\mathbf1-\frac{\ii}{\hbar}\vect{\op p}\cdot\dd\vect x\right)
=\op x+\frac{\ii}{\hbar}[\vect{\op p}\cdot\dd\vect x,\op x]
=\op x+\dd\vect x.\]

\subsection*{2) Schr\"odinger and Heisenberg pictures of time evolution}
In the \emph{Schr\"odinger picture}, the state evolves as
\[\ket{\alpha,t}_{\mathrm S}=\op U(t)\ket{\alpha,0}_{\mathrm S},\qquad
\op U(t)=e^{-\ii\op Ht/\hbar},\]
while the observable $\op A_{\mathrm S}$ remains fixed in time.

In the \emph{Heisenberg picture}, the observable evolves as
\[\op A_{\mathrm H}(t)=\op U^{\dagger}(t)\op A_{\mathrm H}(0)\op U(t),\]
while the state $\ket\alpha_{\mathrm H}$ remains fixed.

For the same initial conditions,
\[\op A_{\mathrm H}(0)=\op A_{\mathrm S},\qquad
\ket{\alpha,0}_{\mathrm S}=\ket\alpha_{\mathrm H}.\]
These are the relations between the two pictures.

\textbf{Notes.}
\begin{enumerate}
\item Inner products (including matrix elements, probability amplitudes and expectation values) are independent of picture:
\[{}_{\mathrm S}\braket{\alpha,t}{\beta,t}_{\mathrm S}
={}_{\mathrm S}\bra{\alpha,0}\op U^{\dagger}\op U\ket{\beta,0}_{\mathrm S}
={}_{\mathrm S}\braket{\alpha,0}{\beta,0}_{\mathrm S}
={}_{\mathrm H}\braket\alpha\beta_{\mathrm H}.\]
\item Since a unitary transformation preserves commutators, if the Schr\"odinger-picture commutation relation is
\[[\op A_{\mathrm S},\op B_{\mathrm S}]=\op C_{\mathrm S},\]
or equivalently initially $[\op A_{\mathrm H}(0),\op B_{\mathrm H}(0)]=\op C_{\mathrm H}(0)$, then the equal-time relation is
\[[\op A_{\mathrm H}(t),\op B_{\mathrm H}(t)]=\op C_{\mathrm H}(t).\]
\item The Hamiltonian in the previously introduced time-evolution operator $\op U(t)=e^{-\ii\op Ht/\hbar}$ is the Schr\"odinger-picture Hamiltonian $\op H=\op H_{\mathrm S}$. For a time-independent Hamiltonian,
\[\op H_{\mathrm H}(t)=\op U^{\dagger}(t)\op H_{\mathrm S}\op U(t)
=e^{\ii\op H_{\mathrm S}t/\hbar}\op H_{\mathrm S}e^{-\ii\op H_{\mathrm S}t/\hbar}
=\op H_{\mathrm S}=\op H_{\mathrm H}.\]
\end{enumerate}
So far we have considered time evolution of states and observables. Do basis vectors also change with time, as general vectors do? Basis vectors are eigenstates of observables; their time dependence is determined by that of the observables.

In the Schr\"odinger picture,
\[\op A_{\mathrm S}\ket a_{\mathrm S}=a\ket a_{\mathrm S}.\]
Since $\op A_{\mathrm S}$ is time independent, its eigenbasis $\ket a_{\mathrm S}$ is also time independent.

In the Heisenberg picture,
\[\op A_{\mathrm H}(t)\ket{a,t}_{\mathrm H}=a\ket{a,t}_{\mathrm H}.\]
Since $\op A_{\mathrm H}$ depends on time, its basis $\ket{a,t}_{\mathrm H}$ also depends on time. Using
\[\op U^{\dagger}(t)\op A_{\mathrm S}\op U(t)\ket{a,t}_{\mathrm H}
=a\ket{a,t}_{\mathrm H},\]
we obtain
\[\op A_{\mathrm S}\op U(t)\ket{a,t}_{\mathrm H}
=a\op U(t)\ket{a,t}_{\mathrm H}.\]
Comparison with the Schr\"odinger-picture eigenvalue equation gives
\[\op U(t)\ket{a,t}_{\mathrm H}=\ket a_{\mathrm S},\quad
\ket{a,t}_{\mathrm H}=\op U^{\dagger}(t)\ket a_{\mathrm S}
=\op U^{\dagger}(t)\ket{a,0}_{\mathrm H}.\]
In summary:
\begin{center}\small
\begin{tabular}{@{}lll@{}}\toprule
&Schr\"odinger picture&Heisenberg picture\\\midrule
State&$\ket{\alpha,t}_{\mathrm S}=\op U(t)\ket{\alpha,0}_{\mathrm S}$&$\ket\alpha_{\mathrm H}$ fixed\\
Observable&$\op A_{\mathrm S}$ fixed&$\op A_{\mathrm H}(t)=\op U^{\dagger}\op A_{\mathrm H}(0)\op U$\\
Basis&$\ket a_{\mathrm S}$ fixed&$\ket{a,t}_{\mathrm H}=\op U^{\dagger}(t)\ket{a,0}_{\mathrm H}$\\\bottomrule
\end{tabular}\end{center}

\subsection*{3) Equation of motion in the Heisenberg picture}
In classical mechanics, only physical quantities evolve with time and the notion of a state vector is not introduced. The Heisenberg picture of quantum mechanics is therefore easier to compare with classical mechanics.

Differentiate
\[\op A_{\mathrm H}(t)=e^{\ii\op Ht/\hbar}\op A_{\mathrm H}(0)e^{-\ii\op Ht/\hbar}\]
with respect to time. We obtain the Heisenberg equation of motion:
\[\boxed{\frac{\dd\op A_{\mathrm H}(t)}{\dd t}
=\frac1{\ii\hbar}[\op A_{\mathrm H}(t),\op H].}\]
It plays a role similar to the Schr\"odinger-picture equation for the state,
\[\ii\hbar\frac{\partial}{\partial t}\ket{\alpha,t}_{\mathrm S}
=\op H\ket{\alpha,t}_{\mathrm S}.\]
Compare the equation of motion of analytical classical mechanics,
\[\frac{\dd A}{\dd t}=\{A,H\}_{\mathrm{Poisson}}.\]
We obtain the route from classical mechanics to quantum mechanics, canonical quantization:
\[\boxed{\{\ \,,\ \}_{\mathrm{Poisson}}\quad
\xrightarrow{\ \text{quantum mechanics}\ }\quad
\frac1{\ii\hbar}[\ \,,\ \ ].}\]

\subsection*{4) Ehrenfest theorem}
How is $\op H$ constructed? For a quantum system with a classical counterpart, the Hamiltonian $H(\vect x,\vect p)$ of analytical mechanics is taken to the quantum Hamiltonian $\op H(\vect{\op x},\vect{\op p})$.

For a free particle,
\[\op H=\frac{\vect{\op p}^{\,2}}{2m}.\]
Heisenberg equations of motion (we henceforth omit the picture index H on observables) give
\[\frac{\dd\vect{\op p}}{\dd t}=
\frac1{\ii\hbar}[\vect{\op p},\op H]=0,\qquad
\vect{\op p}(t)=\vect{\op p}(0),\]
\[\frac{\dd\vect{\op x}}{\dd t}=\frac1{\ii\hbar}[\vect{\op x},\op H]
=\frac1{2\ii\hbar m}[\vect{\op x},\vect{\op p}^{\,2}]
=\frac{\vect{\op p}(t)}m=\frac{\vect{\op p}(0)}m,\quad
\vect{\op x}(t)=\vect{\op x}(0)+\frac{\vect{\op p}(0)}m t.\]
For a general system,
\[\op H=\frac{\vect{\op p}^{\,2}}{2m}+V(\vect{\op x}).\]
The Heisenberg equations are
\[\frac{\dd\vect{\op p}}{\dd t}=
\frac1{\ii\hbar}[\vect{\op p},\op H]
=\frac1{\ii\hbar}[\vect{\op p},V(\vect{\op x})].\]
Expand $V(\vect{\op x})$ as a power series in $\vect{\op x}$ and use $[\op x_i,\op p_j]=\ii\hbar\delta_{ij}$. Then
\[\frac{\dd\vect{\op p}}{\dd t}=-\vect\nabla V(\vect{\op x}),\]
\[\frac{\dd\vect{\op x}}{\dd t}
=\frac1{\ii\hbar}[\vect{\op x},\op H]
=\frac1{2\ii\hbar m}[\vect{\op x},\vect{\op p}^{\,2}]
=\frac{\vect{\op p}}m,\]
\[m\frac{\dd^2\vect{\op x}}{\dd t^2}
=\frac{\dd\vect{\op p}}{\dd t}=-\vect\nabla V(\vect{\op x}).\]
Taking expectation values on both sides,
\[\boxed{m\frac{\dd^2\langle\vect{\op x}\rangle}{\dd t^2}
=\frac{\dd\langle\vect{\op p}\rangle}{\dd t}
=-\langle\vect\nabla V(\vect{\op x})\rangle.}\]
This is the Ehrenfest theorem, corresponding to Newton's law in classical mechanics. Since inner products are picture-independent, the Ehrenfest theorem for expectation values is also picture-independent.

Unless otherwise stated, problems below are generally discussed in the Schr\"odinger picture.